% 第 6 篇（Feng–Harris）末片：印刷页 555–568（PDF 568–581）
% 含附录 A.14–A.19 的余部、参考文献与两位作者通讯地址；不设 \chapter。
% 说明：底本正文实际止于印刷页 568（含参考文献与地址），任务书记为 555–567（13 页）
% 少算一页，故本片给出 555–568 共 14 个页锚。
%
% 与上片衔接：上片 ch-a6-07.tex 止于 PDF 567（印刷 554），末为 A.14 小节与定义 A.14.1
% （André–Quillen 同调群 D_i(B/A;M) 与上同调群 D^i(B/A;M)）。本片自 PDF 568 页首
% （“As a consequence of the properties of the cotangent complex…”）起，不接句子。
\ysrc{555}

作为余切复形性质的推论，André--Quillen（上）同调具有如下若干性质：

\begin{itemize}
\item 设 $B$ 是 $A$-代数，$M'\to M\to M''$ 是 $B$-模的短正合列，则存在长正合列
\[
D^i(B/A;M')\to D^i(B/A;M)\to D^i(B/A;M'')\to D^{i+1}(B/A;M')\to\cdots
\]
\item 设 $B,B'$ 是 $A$-代数，且对 $i>0$ 有 $\Tor_i^A(B,B')=0$。则对任意 $B\otimes_A B'$-模 $M$，有
\[
D^i(B\otimes_A B'/B';M)\cong D^i(B/A;M)
\]
以及
\[
D^i(B\otimes_A B'/A;M)\cong D^i(B/A;M)\oplus D^i(B'/A;M).
\]
\item 若 $A\to B\to C$ 是一列环同态，$M$ 是 $C$-模，则得到长正合列
\[
D^i(B/A;M)\to D^i(C/A;M)\to D^i(C/B;M)\to D^{i+1}(B/A;M)\to\cdots
\]
\end{itemize}

下面我们用 André--Quillen 上同调研究形变理论。

\subsubsection*{A.14.1.\quad 形变理论的设置}

设 $f\colon X\to S$ 是概形间的态射\tn{底本此处作「Let $f\colon X\to S$ be a scheme」，疑脱 morphism of（原书 p.555）。}，$S\hookrightarrow S'$ 是理想层为 $\mathcal I$ 的平方零加厚。我们考虑把 $X$ 形变到 $X'\to S'$，如下图。

\[
\begin{tikzcd}
X \arrow[r, dashed] \arrow[d] & X' \arrow[d, dashed] \\
S \arrow[r, hook] & S'
\end{tikzcd}
\]

于是问题归结为构造扩张

\[
\begin{tikzcd}
0 \arrow[r] & \mathcal J \arrow[r, dashed] & \OO_{X'} \arrow[r, dashed] & \OO_X \arrow[r] & 0 \\
0 \arrow[r] & \mathcal I \arrow[r] \arrow[u, dashed] & \OO_{S'} \arrow[r] \arrow[u, dashed] & \OO_S \arrow[r] \arrow[u] & 0
\end{tikzcd}
\]

其中由于平方零性质，$\mathcal J$ 可看作 $\OO_X$ 上的层；这样看时，它同构于 $f^*\mathcal I\cong\mathcal I\otimes_{\OO_S}\OO_X$。

\subsubsection*{A.14.2.\quad 平方零代数扩张}

受此启发，我们考虑如下问题。设 $f\colon A\to B$ 是（经典）交换环的同态，$M$ 是 $B$-模。用 $M$ 对 $B$ 的\emph{$A$-代数扩张}，指的是正合列
\[
0\to M\to E\to B\to 0,
\]
它把 $E$ 呈现为用 $M$ 对 $B$ 的平方零 $A$-代数扩张。我们希望理解所有这类扩张的结构。结果表明，起控制作用的是 André--Quillen 同调。

\begin{prop}{A.14.2}\label{thm:a6-08-1}
设 $M$ 是 $B$-模。则 $D^1(B/A;M)$ 是用 $M$ 对 $B$ 的 $A$-代数扩张的同构类集合。
\end{prop}

\begin{exam}{A.14.3}\label{thm:a6-08-2}
零元 $0\in D^1(B/A;M)$ 对应平凡平方零扩张 $B\oplus M$，其乘法为
\[
(b_1,m_1)(b_2,m_2)=(b_1b_2,\ b_1m_2+b_2m_1).
\]
\end{exam}

\begin{proof}
我们给出两个方向的映射。先设给定扩张 $M\to E\to B$。取自由 $A$-代数消解 $A\to\mathcal B\xrightarrow{\sim}B$。特别地，由于 $\mathcal B_0$ 是 $A$ 上的多项式代数，存在提升
\[
\begin{tikzcd}
\mathcal B_0 \arrow[r] \arrow[d, dashed, "\phi_0"] & B \arrow[d, equal] \\
M \arrow[r] & E \arrow[r] & B
\end{tikzcd}
\]
现在有两个映射 $d_0,d_1\colon\mathcal B_1\rightrightarrows\mathcal B_0$，且已知 $(d_0-d_1)$ 的像在 $\mathcal B_0$ 中生成理想，并有 $\mathcal B_0/\operatorname{Im}(d_0-d_1)\cong\pi_0(\mathcal B)\cong B$。于是差
\[
\phi_0\circ(d_0-d_1)\colon\mathcal B_1\to M
\]
是从 $\mathcal B_1$ 到 $M$ 的 $B$-线性导子，因而对应 $\Hom_B(\Omega_{\mathcal B_1},M)$ 中的一个元素。可以验证它是上闭链，故诱导 $D^1(B/A;M)$ 中的一个类。$\phi_0$ 的不同选取只把这上闭链改变一个上边缘，因此从扩张的同构类到 $D^1(B/A;M)$ 的映射是良定的。

这同时提示如何构造逆映射。给定 $D^1(B/A;M)$ 中由导子 $\delta\colon\mathcal B_1\to M$ 表示的类，定义 $E$ 为 $B$-模的推出。
\[
\begin{tikzcd}
\mathcal B_1 \arrow[r, "d_0-d_1"] \arrow[d, "\delta"] & \mathcal B_0 \arrow[d, dashed] \\
M \arrow[r] & E
\end{tikzcd}
\]
再给 $E$ 配上由它作为 $\mathcal B_0\oplus M$ 的商而诱导的乘法。为使其良定，只需验证 $\mathcal B_1$ 在 $(d_1-d_0,\delta)$ 下的像是平方零理想，而这是直接的。
\end{proof}

\begin{exam}{A.14.4}\label{thm:a6-08-3}
考虑用 $\F_p$ 对 $\F_p$ 的 $\Z$-代数扩张。由于我们容许交换环的任意扩张，这类扩张由 $D^1(\F_p/\Z;\F_p)$ 控制。由命题 A.12.2，$\Z\to\F_p$ 的余切复形为 $\mathbf L_{\F_p/\Z}\cong\F_p[1]$。故
\[
D^1(\F_p/\Z;\F_p)\cong\Hom_{\F_p}(\F_p[1],\F_p[1])\cong\F_p.
\]
0 类对应扩张 $\F_p[\epsilon]/\epsilon^2$。非零类的底环都是 $\Z/p^2\Z$，到 $\F_p$ 的映射是自然投影与 $\F_p$ 的某个自同构的复合。
\end{exam}

\begin{exam}{A.14.5}\label{thm:a6-08-4}
设 $A\to B$ 是交换环的光滑同态，$M$ 是任意 $B$-模。则由命题 A.12.1，$\mathbf L_{B/A}\cong\Omega^1_{B/A}$ 有限投射，故 $D^1(B/A;M)=0$。这说明用 $M$ 对 $B$ 的平方零 $A$-代数扩张唯一，且必为分裂扩张。
\end{exam}

\ysrc{556}

\begin{proof}[命题 A.12.1 证明的补全]
我们想证明：若 $f\colon A\to B$ 是 Noether 环的有限型态射，且 $\mathbf L_{B/A}$ 是集中在 0 次的有限投射 $B$-模，则 $f$ 光滑。借助有限性假设，只需证明 $f$ 形式光滑，即

\ysrc{557}

对任意平方零扩张 $S\hookrightarrow\tilde S$，图表
\[
\begin{tikzcd}
S \arrow[r] \arrow[d] & \Spec B \arrow[d] \\
\tilde S \arrow[r] \arrow[ur, dashed] & \Spec A
\end{tikzcd}
\]
都有提升。先设映射 $S\to\Spec B$ 是同构。此时 $\tilde S$ 是某个用 $M$ 对 $B$ 的平方零 $A$-代数扩张的谱，我们想说明它分裂。但关于 $\mathbf L_{B/A}$ 的假设蕴含：对任意 $B$-模 $M$ 都有 $D^1(B/A;M)=0$。于是 $[\tilde S]\in D^1(B/A;M)$ 的类为零，这恰好说明存在所需的提升。

现处理一般情形。设 $S=\Spec R$，$\tilde S=\Spec\tilde R$，并令 $I=\ker(\tilde R\to R)$。取从多项式 $A$-代数 $P$ 到 $R$ 的满射，设其核为 $J$。通过选取生成元的提升，我们得到提升 $P\to\tilde R$；它把 $J$ 送入 $I$，因而经过 $P/J^2$ 分解。于是 $P/J^2$ 是用 $M$ 对 $B$ 的平方零 $A$-代数扩张，由上一段处理的情形存在分裂 $B\to P/J^2$；把它与到 $\tilde R$ 的映射复合，即得所需提升。
\[
\begin{tikzcd}[row sep=large]
P \arrow[r, "\sim"] \arrow[d, equal] & B=P/J \arrow[r] \arrow[d, dashed, bend left=55] & R=\tilde R/I \arrow[d] \\
P \arrow[r] & P/J^2 \arrow[u] & \tilde R
\end{tikzcd}
\]
\end{proof}

\subsection*{A.15.\quad 概形的形变理论}

设 $\tilde A\to A$ 是用平方零理想 $I$ 所作的扩张。我们考虑寻找平坦形变的问题
\[
\begin{tikzcd}
J \arrow[r] & \tilde B \arrow[r] & B \\
I \arrow[r] \arrow[u] & \tilde A \arrow[r] \arrow[u] & A \arrow[u]
\end{tikzcd}
\]
或者用几何语言：把 $\Spec B\to\Spec A$ 形变为 $\Spec\tilde A$ 上的平坦族。
\[
\begin{tikzcd}
\Spec B \arrow[d] & \Spec\tilde B \arrow[l, hook] \arrow[d] \\
\Spec A & \Spec\tilde A \arrow[l, hook]
\end{tikzcd}
\]
由于 $I$ 平方零，我们可把 $I$ 看作 $A$-模。同样，可把 $J$ 看作 $B$-模。

\begin{lem}{A.15.1}\label{thm:a6-08-5}
在上述情形中，$\tilde B$ 在 $\tilde A$ 上平坦，当且仅当 $I\otimes_A B\xrightarrow{\sim}J$。
\end{lem}

\begin{proof}
若 $\tilde B$ 在 $\tilde A$ 上平坦，则把 $\tilde B\otimes_{\tilde A}(-)$ 作用于短正合列 $I\to\tilde A\to A$，并利用 $I$ 上的 $\tilde A$-作用经过 $A$ 分解，即得 $I\otimes_A B\xrightarrow{\sim}J$。
另一方向见 [Sta22, Tag 00MD]。
\end{proof}

\ysrc{558}

今后我们可假设 $J=I\otimes_A B$。已知用 $J$ 对 $B$ 的 $\tilde A$-代数扩张由 $D^1(B/\tilde A;J)$ 分类。复合 $\tilde A\to A\to B$ 给出长正合列
\[
\begin{aligned}
&\to D^1(B/A;J)\to D^1(B/\tilde A;J)\to D^1(A/\tilde A;J)\\
&\to D^2(B/A;J)
\end{aligned}
\]
扩张 $\tilde A$ 的等价类可视为元素 $[\tilde A]\in D^1(A/\tilde A;J)$；它来自 $D^1(B/\tilde A;J)$ 当且仅当它在 $D^2(B/A;J)$ 中的像为零。由长正合列，它在 $D^1(B/\tilde A;J)$ 中的原像构成 $D^1(B/A;J)$ 上的挠子（可能为空）。于是得到：

\begin{prop}{A.15.2}\label{thm:a6-08-6}
设情形如上。
\begin{itemize}
\item 与扩张 $\tilde A\to A$ 相伴的有一个阻碍类 $\mathrm{obs}(\tilde A)\in D^2(B/A;J)$，它为零当且仅当存在扩张 $\tilde B$。
\item 若 $\mathrm{obs}(\tilde A)\in D^2(B/A;J)$ 为零，则扩张的同构类构成 $D^1(B/A;J)$ 上的非空挠子。
\item 任意扩张的自同构群是
\[
D^0(B/A;J)\cong\Hom_B(\Omega_{B/A},J).
\]
\end{itemize}
（这是初等的；与余切复形理论无关。）
\end{prop}

\subsubsection*{A.15.1.\quad 应用：Witt 向量}

下面把前述理论应用于完美 $\F_p$-代数的 Witt 向量的构造。

\begin{lem}{A.15.3}\label{thm:a6-08-7}
若 $B$ 是完美 $\F_p$-代数，则 $\mathbf L_{B/\F_p}=0$。
\end{lem}

\begin{proof}
按定义，$B$ 上的 Frobenius 自同态是同构，因而诱导 $\mathbf L_{B/\F_p}$ 的自同构。为计算 Frobenius 在 $\mathbf L_{B/\F_p}$ 上的作用，取 $\F_p$ 上自由消解 $\mathcal B\xrightarrow{\sim}B$，并注意 $B$ 上的 $\Frob$ 可提升为 $\mathcal B$ 上的 $\Frob$（在 $\mathcal B$ 上逐层取 $p$ 次幂映射）：
\[
\begin{tikzcd}
\mathcal B \arrow[r, "\Frob"] \arrow[d, "\sim"] & \mathcal B \arrow[d, "\sim"] \\
B \arrow[r, "\Frob"] & B
\end{tikzcd}
\]
因此 $\Frob$ 在 $\mathbf L_{B/\F_p}$ 上的作用，是由它在 $\mathbf L_{\mathcal B/\F_p}$ 上的作用经张量得到的；而后者为零，因为 Frobenius 自同态在 Kähler 微分上诱导 0。于是一方面 $\Frob$ 诱导 $\mathbf L_{B/\F_p}$ 的自同构，另一方面它又是零映射。故 $\mathbf L_{B/\F_p}=0$。
\end{proof}

\begin{lem}{A.15.4}\label{thm:a6-08-8}
$B$ 到 $\Z/p^n$ 的平坦形变唯一。
\end{lem}

\begin{proof}
$n=1$ 的基情形是同义反复。用归纳法，设已得到 $A_n:=\Z/p^n$ 上的唯一形变 $B_n$，并希望把它延拓到 $\tilde A=\Z/p^{n+1}$ 上。令 $I_n\subset\tilde A$ 为 $\tilde A\to A_n$ 的核；它是 $A_n$-模，事实上由 $A_1$ 拉回得到。令 $J_n=I_n\otimes_{A_n}B_n$，它由 $B$ 拉回得到。由于 $B_n$ 在 $J_n$ 上的作用经过 $B_1=B$ 分解，且 $B_n$ 在 $A_n$ 上平坦，对每个 $i$ 有
\[
D^i(B_n/A_n;J_n)\cong D^i(B/A_1;J_n)\cong 0.
\]
特别地，命题 A.15.2 中的阻碍类落在 $D^2(B_n/A_n;J_n)=0$ 中，故扩张存在；又因 $D^1(B_n/A_n;J_n)=0$，该扩张唯一。
\end{proof}

\ysrc{559}

因此存在唯一的图表（在同构意义下唯一）
\[
\begin{tikzcd}
B & B_2 \arrow[l, hook] & B_3 \arrow[l, hook] & \cdots \arrow[l, hook] \\
\F_p \arrow[u] & \Z/p^2\Z \arrow[u] \arrow[l, hook] & \Z/p^3\Z \arrow[u] \arrow[l, hook] & \cdots \arrow[u] \arrow[l, hook]
\end{tikzcd}
\]
其中所有方块余笛卡儿，所有竖箭头平坦。逆极限 $\varprojlim_n B_n$ 称为 $B$ 的\emph{Witt 向量环}。Witt 向量的这种构造由 [Sch12, Remark 5.14] 指出。

\subsection*{A.16.\quad 映射的形变理论}

下一个情形来自映射的形变理论：给定平方零扩张 $X\hookrightarrow\tilde X$ 与映射 $X\to Y$，我们希望理解映射的扩张
\[
\begin{tikzcd}
X \arrow[d] & \tilde X \arrow[l, hook] \arrow[dl, dashed] \\
Y
\end{tikzcd}
\]
局部问题是：给定 $A$-代数 $B$ 与平方零扩张 $J\to\tilde B\to B$，给 $\tilde B$ 配上相容的 $A$-代数结构。把 $\tilde B$ 看作用 $J$ 对 $B$ 的交换 $\Z$-代数扩张，其等价类可视为元素 $[\tilde B]\in D^1(B/\Z;J)$。我们希望把它提升为 $D^1(B/A;J)$ 中的一个类\tn{底本此处作「We want to lift this a class in $D^1(B/A;J)$」，疑脱 to（原书 p.559）。}。我们有长正合列
\[
\cdots\to D^0(A/\Z;J)\to D^1(B/A;J)\to D^1(B/\Z;J)\to D^1(A/\Z;J)\to\cdots
\]
映射 $D^1(B/\Z;J)\to D^1(A/\Z;J)$ 可解释为把 $\tilde B$ 送到纤维积 $A\times_B\tilde B$，后者是用 $J$ 对 $A$ 的交换代数扩张。于是可把长正合列解释如下。

\begin{prop}{A.16.1}\label{thm:a6-08-9}
设情形如上。
\begin{itemize}
\item 与扩张 $\tilde B\to B$ 相伴的有一个阻碍类 $\mathrm{obs}(\tilde B)\in D^1(A/\Z;J)$，它为零当且仅当 $\tilde B$ 上存在与 $B$ 上给定结构相容的 $A$-代数结构。
\item 若 $\mathrm{obs}(\tilde B)\in D^1(A/\Z;J)$ 为零，则 $A$-代数结构的同构类构成 $D^0(A/\Z;J)=\Hom_A(\Omega_{A/\Z},J)$ 上的非空挠子。
\end{itemize}
（这是初等的；与余切复形理论无关。）
\end{prop}

\subsubsection*{A.16.1.\quad 应用：Fontaine 映射}

设 $R$ 是 $p$-进完备且 $p$-无挠的环。令 $R^\flat:=\varprojlim_{\Frob_p}(R/p)$。这是完美环；作为映射到 $R/p$ 的「终」完美环，它在同构意义下被唯一确定。

$p$-进 Hodge 理论中一个非常重要的事实是：存在唯一的映射 $W(R^\flat)\to R$，它模 $p$ 约化后即典范映射 $R^\flat\to R/p$，称为 \emph{Fontaine 映射}；不过本文无法解释它的意义。这里 $W(R^\flat)$

\ysrc{560}

是 Witt 向量；如 §A.15.1 所述，它是唯一（在同构意义下唯一）的一族提升的逆极限：
\[
\begin{tikzcd}
R^\flat & W_2(R^\flat) \arrow[l, hook] & W_3(R^\flat) \arrow[l, hook] & \cdots \arrow[l, hook] \\
\F_p \arrow[u] & \Z/p^2\Z \arrow[u] \arrow[l, hook] & \Z/p^3\Z \arrow[u] \arrow[l, hook] & \cdots \arrow[u] \arrow[l, hook]
\end{tikzcd}
\]
下面逐步构造一族映射，从 $W_2(R^\flat)\to R/p^2$ 开始。
\[
\begin{tikzcd}[row sep=large]
R/p & R/p^2 \arrow[l, hook] & R/p^3 \arrow[l, hook] & \cdots \arrow[l, hook] \\
R^\flat \arrow[u] & W_2(R^\flat) \arrow[u, dashed] \arrow[l, hook] & W_3(R^\flat) \arrow[u, dashed] \arrow[l, hook] & \cdots \arrow[u] \arrow[l, hook] \\
\F_p \arrow[u] & \Z/p^2\Z \arrow[u] \arrow[l, hook] & \Z/p^3\Z \arrow[u] \arrow[l, hook] & \cdots \arrow[u] \arrow[l, hook]
\end{tikzcd}
\]
（$R$ 的 $p$-无挠性用来保证 $R/p^n$ 在 $\Z/p^n\Z$ 上平坦。）用与上面类似的论证，这类扩张的阻碍是
\[
D^1\bigl(W_2(R^\flat)/(\Z/p^2\Z);\,pW_2(R^\flat)\bigr)
\]
中的一个类，而它同构于 $D^1(R^\flat/\F_p;R^\flat)$；由于 $R^\flat$ 完美，后者为零。于是这类扩张的集合\tn{底本此处作「the set of such an extensions」，an 为衍字（原书 p.560）。}是 $D^0(R^\flat/\F_p;R^\flat)=0$ 上的挠子，故扩张唯一。归纳地，我们对每个 $n\ge1$ 构造了相容的映射 $W_n(R^\flat)\to R/p^n$。最后，取这些映射的逆极限即得 Fontaine 映射（这里用到 $R$ 的 $p$-进完备性）。

\subsection*{A.17.\quad 整体形变}

迄今我们讨论的是仿射概形的形变。现在我们开始研究可能非仿射的概形的形变。

\begin{exam}{A.17.1}\label{thm:a6-08-10}
设 $X$ 是域 $k$ 上的光滑簇（不必仿射）。$X$ 的\emph{一阶形变}指的是 $k[\epsilon]/(\epsilon^2)$ 上的平坦族 $\tilde X$，它在 $\Spec k$ 上的纤维为 $X$。我们断言：$X$ 的一阶形变集合与 $H^1(X,TX)$ 一一对应。

为看清这一点，写 $X=\bigcup U_i$ 为仿射开集 $U_i=\Spec A_i$ 的并。$X$ 的一阶形变由各 $U_i$ 的一阶形变的一组相容数据唯一决定。由于 $U_i$ 光滑，它们有唯一的一阶形变 $\tilde U_i$。粘接数据是 $\tilde U_i\cap\tilde U_j:=\Spec\tilde A_{ij}$ 的一个自同构，它提升 $U_i\cap U_j:=\Spec A_{ij}$ 上给定的自同构，形如
\[
a+b\epsilon\mapsto a+(\delta a+b)\epsilon,
\]
其中 $\delta$ 是从 $A_{ij}$ 到自身的 $k$-线性导子，即切层 $TX$ 在 $U_{ij}$ 上的一个截面。为粘成 $X$ 的一阶形变，这些截面须在三重交上满足上闭链条件。最后，用自同构（作为一阶形变）修改每个 $\tilde U_i$，相当于加上一个上边缘。于是我们看到，一阶形变的数据恰是 Čech 上同调 $H^1(X,TX)$ 中一个类的数据。

作为应用，考虑\emph{（光滑、射影）亏格 $g$ 曲线的模空间} $\mathcal M_g$。该模空间的 $k$-点是光滑射影亏格 $g$ 曲线 $C$。$\mathcal M_g$ 在对应于 $C$ 的 $k$-点处的切空间，其底层集合是\ysrc{561}把该 $k$-点扩张为 $k[\epsilon]/\epsilon^2$-点的扩张集合，也就是 $C/k$ 的一阶形变集合。按上述结论，这就是 $H^1(C,TC)$。此外，阻碍取值于 $H^2(C,TC)$；由于 $C$ 是曲线，它为零——这说明 $\mathcal M_g$ 在对应于 $C$ 的 $k$-点处光滑。
\end{exam}

\subsection*{A.18.\quad 余切复形的整体化}

把命题 A.15.2 整体化，可得：

\begin{prop}{A.18.1}\label{thm:a6-08-11}
设 $f\colon X\to S$ 是概形间的映射。设 $S\hookrightarrow\tilde S$ 是理想层为 $\mathcal I$ 的平方零加厚，$\mathcal J=f^*\mathcal I$。
\begin{itemize}
\item 与扩张 $S\hookrightarrow\tilde S$ 相伴的有一个阻碍类 $\mathrm{obs}(\tilde S)\in\Ext^2(\mathbf L_{X/S},\mathcal J)$，它为零当且仅当存在平坦扩张 $\tilde X\to\tilde S$。
\item 若 $\mathrm{obs}(\tilde S)\in\Ext^2(\mathbf L_{X/S},\mathcal J)$ 为零，则平坦扩张的同构类构成 $\Ext^1(\mathbf L_{X/S},\mathcal J)$ 上的非空挠子。
\item 任意平坦扩张的自同构群是 $\Ext^0(\mathbf L_{X/S},\mathcal J)$。
\end{itemize}
\end{prop}

\begin{exam}{A.18.2（LCI 曲线）}\label{thm:a6-08-12}
回忆：Noether 环的有限型同态 $A\to B$ 称为\emph{局部完全交}（LCI），如果它 Zariski 局部地可写成复合 $A\to A[x_1,\dots,x_n]\twoheadrightarrow B$，其中第二个映射是正则商。更一般地，概形的态射称为 LCI，如果它 Zariski 局部是 LCI。例如，域上至多带节点奇点的曲线是 LCI。由命题 A.12.2，若 $Y\to X$ 是 LCI，则 $\mathbf L_{Y/X}$ 集中在次数 $-1,0$。

设 $C$ 是域 $k$ 上的 LCI 曲线，它在有限个闭点之外光滑\tn{底本此处作「which is is smooth outside」，is 重出（原书 p.561）。}。我们断言 $\Ext^2(\mathbf L_{C/k},\mathcal J)=0$，这特别说明 $C/k$ 的所有无穷小形变都无阻碍。由局部--整体谱序列，只需证明以下三项为零：
\begin{itemize}
\item $H^2(C,\Ext^0(\mathbf L_{C/k},\mathcal J))$，
\item $H^1(C,\Ext^1(\mathbf L_{C/k},\mathcal J))$，以及
\item $H^0(C,\Ext^2(\mathbf L_{C/k},\mathcal J))$。
\end{itemize}
第一项为零，因为 $C$ 是曲线，其上同调维数为 1。第三项为零，因为 $\Ext^2(\mathbf L_{C/k},\mathcal J)$ 为零，这要归功于 $\mathbf L_{C/k}$ 支撑在次数 $[-1,0]$ 中。最后，由于 $C$ 在有限个闭点之外光滑，由余切复形的局部化性质，再加上它在光滑情形的计算（§A.11），$\Ext^1(\mathbf L_{C/k},\mathcal J)$ 是挠层，故其更高上同调为零。
\end{exam}

为完整起见，我们陈述命题 A.16.1 的整体化：

\begin{prop}{A.18.3}\label{thm:a6-08-13}
设 $X\hookrightarrow\tilde X$ 是理想层为 $\mathcal J$ 的平方零扩张，$f\colon X\to Y$。设 $\mathbf L_Y$ 是 $Y$ 的绝对余切复形（即 $Y\to\Spec\Z$ 的余切复形）。
\begin{itemize}
\item 与 $\tilde X$ 相伴的有一个阻碍类 $\mathrm{obs}(\tilde X)\in\Ext^1(f^*\mathbf L_Y,\mathcal J)$，它为零当且仅当 $f$ 可延拓为映射 $\tilde X\to Y$。
\item 若 $\mathrm{obs}(\tilde X)\in\Ext^1(f^*\mathbf L_Y,\mathcal J)$ 为零，则 $f$ 的扩张的同构类构成 $\Ext^0(f^*\mathbf L_Y,\mathcal J)$ 上的非空挠子。
\end{itemize}
\end{prop}

\ysrc{562}

\subsection*{A.19.\quad André--Quillen 同调的几何解释}

本节反思导出代数几何与余切复形之间的关系；更详尽的论述可参看 [Vez10]。正如 §3 所讨论的，导出代数几何在实践中的一个特征性能力，是构造出余切复形可被「有意义的」识别（例如识别为某个自然上同调理论）的空间。之所以如此，是因为导出代数几何为余切复形提供了「几何」解释。

考虑概形与约化概形之间的类比（见 §2.3）或许有帮助。即使我们只关心约化概形，把非约化概形纳入词汇，也让我们能把切空间解释为从对偶数出发的映射\tn{底本此处作「allows us to gives」，gives 疑为 give 之误（原书 p.562）。}。在模空间的语境中，这让我们能把切空间表达为一阶形变，进而在实践中常常表达为一个上同调群。

我们可以设想把「约化模空间」定义为约化交换环上的函子。但先验地看，如何显式描述它的切空间并不清楚。若该函子能延拓到非约化环，则其切空间可按上一段计算；又若该模空间恰好光滑，人们就会后验地知道底层约化概形的切空间。读者不妨在 $\mathcal M_g$（光滑射影亏格 $g$ 曲线的模空间）这个例子上思考这一思想实验。

同样地，经典模空间的\emph{切复形}先验地难以显式描述，因为「更高」的切空间并不由某个从模问题出发的直接过程给出。然而在导出世界里，完整的切复形确实容许这样的直接描述，我们下面来说明。

设 $A\to B$ 是交换环的同态，$\eta\colon B\to A$ 是增广。聚焦于特殊情形 $A=\Spec k$（$k$ 为域）或许有帮助，此时 $\eta$ 对应 $B$ 的一个 $k$-点。

设 $M$ 是 $B$-模。它可变成单纯 $B$-模 $M[0]$，其特征性质是：它是唯一满足 $\pi_0(M[0])=M$ 且 $i>0$ 时 $\pi_i(M[0])=0$ 的单纯 $B$-模。对 $n\ge0$，存在单纯 $B$-模 $M[n]$，满足 $\pi_n(M[n])=M$，且 $i\ne n$ 时 $\pi_i(M[n])=0$。（例如可用 Dold--Kan 对应来构造 $M[n]$。）

然后我们造出逐层的单纯平方零扩张 $A\oplus M[n]$。其分次同伦环是用 $M$（在次数 $n$）对 $A$（在次数 0）的平方零扩张。

\begin{exam}{A.19.1}\label{thm:a6-08-14}
设 $A=k$ 是域，$M=k$。则 $k\oplus k[0]=k[\epsilon]/\epsilon^2$ 是 $k$ 上的\emph{对偶数环}。更一般地，$k\oplus k[n]$ 称为 $k$ 上的\emph{$n$ 阶导出对偶数环}。
\end{exam}

于是 André--Quillen 上同调群 $D^n(B/A;M)$ 可看作从 $B$ 到 $A\oplus M[n]$ 且提升 $\eta$ 的同态群，对象是单纯 $A$-代数的同伦范畴。理由相当形式化：只要相关定义就位即可；这里只给一个梗概。为计算这些同态，按定义需用自由的单纯 $A$-代数 $\mathcal B$ 消解 $B$。（这类似于导出 $A$-模范畴中的映射：\ysrc{563}先要用投射消解来消解源对象。）然后，提升
\[
\begin{tikzcd}
& A\oplus M[n] \\
\mathcal B \arrow[r] \arrow[ur, dashed] & A \arrow[u]
\end{tikzcd}
\]
由从 $\mathcal B$ 到 $M[n]$ 的导子计算；后者又可转述为从 $\mathcal B$ 的 Kähler 微分到 $M[n]$ 的同态，这基本上就是 André--Quillen 上同调的定义。

几何上，我们可把「$\Spec(A\oplus M[n])$」设想为用 $M$ 对 $A$ 的某种「更高导出无穷小加厚」。导出模问题可在这类空间上取值，从而给出 André--Quillen 同调群的几何解释；这使得导出模空间的余切复形比经典模空间的余切复形更容易理解，其意义与「经典模空间的切空间比『约化模空间』的切空间更容易理解」相同。

\section*{参考文献}

\begin{refs}
\item[{[AHHV17]}] N. Abe, G. Henniart, F. Herzig, and M.-F. Vign\'eras, \emph{A classification of irreducible admissible mod $p$ representations of $p$-adic reductive groups}, J. Amer. Math. Soc. \textbf{30} (2017), no. 2, 495--559, DOI 10.1090/jams/862. MR3600042
\item[{[AM86]}] M. Artin and B. Mazur, \emph{Etale homotopy}, reprint of the 1969 original, Lecture Notes in Mathematics, vol. 100, Springer-Verlag, Berlin, 1986. MR883959
\item[{[Arb21]}] \emph{Arbeitsgemeinschaft: Derived Galois Deformation Rings and Cohomology of Arithmetic Groups (hybrid meeting)}, Oberwolfach Rep. \textbf{18} (2021), no. 2, 1001--1046, DOI 10.4171/owr/2021/18. Abstracts from the working session held April 4--10, 2021; Organized by Frank Calegari, S\o ren Galatius and Akshay Venkatesh. MR4472298
\item[{[Ata22]}] Stanislav Ivanov Atanasov, \emph{Derived Hecke Operators on Unitary Shimura Varieties}, ProQuest LLC, Ann Arbor, MI, 2022. Thesis (Ph.D.)--Columbia University. MR4435548
\item[{[Avr99]}] Luchezar L. Avramov, \emph{Locally complete intersection homomorphisms and a conjecture of Quillen on the vanishing of cotangent homology}, Ann. of Math. (2) \textbf{150} (1999), no. 2, 455--487, DOI 10.2307/121087. MR1726700
\item[{[BDSBT01]}] Kevin Buzzard, Mark Dickinson, Nick Shepherd-Barron, and Richard Taylor, \emph{On icosahedral Artin representations}, Duke Math. J. \textbf{109} (2001), no. 2, 283--318, DOI 10.1215/S0012-7094-01-10922-8. MR1845181
\item[{[Ber11]}] Laurent Berger, \emph{Trianguline representations}, Bull. Lond. Math. Soc. \textbf{43} (2011), no. 4, 619--635, DOI 10.1112/blms/bdr036. MR2820149
\item[{[BHR94]}] Don Blasius, Michael Harris, and Dinakar Ramakrishnan, \emph{Coherent cohomology, limits of discrete series, and Galois conjugation}, Duke Math. J. \textbf{73} (1994), no. 3, 647--685, DOI 10.1215/S0012-7094-94-07326-2. MR1262930
\item[{[BIP23]}] Gebhard B\"ockle, Ashwin Iyengar, and Vytautas Pa\v{s}k\={u}nas, \emph{On local Galois deformation rings}, Forum Math. Pi \textbf{11} (2023), Paper No. e30, 54, DOI 10.1017/fmp.2023.25. MR4660960
\item[{[BKM21]}] Gebhard B\"ockle, Chandrashekhar B. Khare, and Jeffrey Manning, \emph{Wiles defect for Hecke algebras that are not complete intersections}, Compos. Math. \textbf{157} (2021), no. 9, 2046--2088, DOI 10.1112/S0010437X21007454. MR4301563
\item[{[BMS19]}] Bhargav Bhatt, Matthew Morrow, and Peter Scholze, \emph{Topological Hochschild homology and integral $p$-adic Hodge theory}, Publ. Math. Inst. Hautes \'Etudes Sci. \textbf{129} (2019), 199--310, DOI 10.1007/s10240-019-00106-9. MR3949030
\item[{[Bor76]}] Armand Borel, \emph{Admissible representations of a semi-simple group over a local field with vectors fixed under an Iwahori subgroup}, Invent. Math. \textbf{35} (1976), 233--259, DOI 10.1007/BF01390139. MR444849

\ysrc{564}
\item[{[BP12]}] Christophe Breuil and Vytautas Pa\v{s}k\={u}nas, \emph{Towards a modulo $p$ Langlands correspondence for $\mathrm{GL}_2$}, Mem. Amer. Math. Soc. \textbf{216} (2012), no. 1016, vi+114, DOI 10.1090/S0065-9266-2011-00623-4. MR2931521
\item[{[BR24]}] Roman Bezrukavnikov and Simon Riche, \emph{On two modular geometric realizations of an affine Hecke algebra}, arXiv:2402:08281.\tn{底本此处作 \texttt{arXiv:2402:08281}，编号中的冒号疑为句点之误（原书 p.564）。}
\item[{[Bre03]}] Christophe Breuil, \emph{Sur quelques repr\'esentations modulaires et $p$-adiques de $\mathrm{GL}_2(\Q_p)$. II} (French, with French summary), J. Inst. Math. Jussieu \textbf{2} (2003), no. 1, 23--58, DOI 10.1017/S1474748003000021. MR1955206
\item[{[BS22]}] Bhargav Bhatt and Peter Scholze, \emph{Prisms and prismatic cohomology}, Ann. of Math. (2) \textbf{196} (2022), no. 3, 1135--1275, DOI 10.4007/annals.2022.196.3.5. MR4502597
\item[{[BS23]}] Bhargav Bhatt and Peter Scholze, \emph{Prismatic $F$-crystals and crystalline Galois representations}, Camb. J. Math. \textbf{11} (2023), no. 2, 507--562.
\item[{[BW00]}] A. Borel and N. Wallach, \emph{Continuous cohomology, discrete subgroups, and representations of reductive groups}, 2nd ed., Mathematical Surveys and Monographs, vol. 67, American Mathematical Society, Providence, RI, 2000, DOI 10.1090/surv/067. MR1721403
\item[{[BZN18]}] David Ben-Zvi and David Nadler, \emph{Betti geometric Langlands}, Algebraic geometry: Salt Lake City 2015, Proc. Sympos. Pure Math., vol. 97, Amer. Math. Soc., Providence, RI, 2018, pp. 3--41, DOI 10.1090/pspum/097.2/01. MR3821166
\item[{[BZSV24]}] David Ben-Zvi, Yiannis Sakellaridis, and Akshay Venkatesh, \emph{Relative Langlands duality}, arXiv:2408.04677, 2024.
\item[{[Cal20]}] Frank Calegari, \emph{Motives and $L$-functions}, Current developments in mathematics 2018, Int. Press, Somerville, MA, 2020, pp. 57--123. MR4363376
\item[{[CDT99]}] Brian Conrad, Fred Diamond, and Richard Taylor, \emph{Modularity of certain potentially Barsotti-Tate Galois representations}, J. Amer. Math. Soc. \textbf{12} (1999), no. 2, 521--567, DOI 10.1090/S0894-0347-99-00287-8. MR1639612
\item[{[CG18]}] Frank Calegari and David Geraghty, \emph{Modularity lifting beyond the Taylor-Wiles method}, Invent. Math. \textbf{211} (2018), no. 1, 297--433, DOI 10.1007/s00222-017-0749-x. MR3742760
\item[{[CS25]}] Ana Caraiani and Sug Woo Shin, \emph{Recent progress on Langlands reciprocity for general linear groups: Shimura varieties and beyond}, The Langlands program: I. Representations, Shimura varieties, and shtukas, Proc. Sympos. Pure Math., vol. 112.1, Amer. Math. Soc., Providence, RI, 2025.
\item[{[DHKM20]}] Jean-Fran\c{c}ois Dat, David Helm, Robert Kurinczuk, and Gilbert Moss, \emph{Moduli of Langlands parameters}, J. Eur. Math. Soc. (JEMS) \textbf{27} (2025), no. 5, 1827--1927, DOI 10.4171/jems/1599. MR4889237
\item[{[DHRV22]}] Henri Darmon, Michael Harris, Victor Rotger, and Akshay Venkatesh, \emph{The derived Hecke algebra for dihedral weight one forms}, Michigan Math. J. \textbf{72} (2022), 145--207, DOI 10.1307/mmj/20217221. MR4460251
\item[{[dJ01]}] A. J. de Jong, \emph{A conjecture on arithmetic fundamental groups}, Israel J. Math. \textbf{121} (2001), 61--84, DOI 10.1007/BF02802496. MR1818381
\item[{[DS74]}] Pierre Deligne and Jean-Pierre Serre, \emph{Formes modulaires de poids 1} (French), Ann. Sci. \'Ecole Norm. Sup. (4) \textbf{7} (1974), 507--530 (1975). MR379379
\item[{[DS95]}] W. G. Dwyer and J. Spali\'nski, \emph{Homotopy theories and model categories}, Handbook of algebraic topology, North-Holland, Amsterdam, 1995, pp. 73--126, DOI 10.1016/B978-044481779-2/50003-1. MR1361887
\item[{[EGH25]}] Matthew Emerton, Toby Gee, and Eugen Hellmann, \emph{An introduction to the categorical $p$-adic Langlands program}, The Langlands program: II. Geometrization of the Langlands correspondence, Proc. Sympos. Pure Math., vol. 112.2, Amer. Math. Soc., Providence, RI, 2025.
\item[{[Fen20]}] Tony Feng, \emph{The spectral Hecke algebra}, arXiv:1912.04413.
\item[{[Fen24]}] Tony Feng, \emph{Smith theory and cyclic base change functoriality}, Forum Math. Pi \textbf{12} (2024), Paper No. e1, 66, DOI 10.1017/fmp.2023.32. MR4689770
\item[{[FG]}] Tony Feng and Dennis Gaitsgory, \emph{Derived Satake isomorphisms and derived commuting varieties}.
\item[{[FHMR]}] Tony Feng, Michael Harris, Barry Mazur, and Arpon Raksit, \emph{Derived Langlands for $\mathrm{GL}(1)$}, In preparation.

\ysrc{565}
\item[{[FK24]}] Tony Feng and Adeel Khan, \emph{Modularity of higher theta series II: Chow group of the generic fiber}, arXiv:2403.19711, 2024.
\item[{[Fra98]}] Jens Franke, \emph{Harmonic analysis in weighted $L^2$-spaces} (English, with English and French summaries), Ann. Sci. \'Ecole Norm. Sup. (4) \textbf{31} (1998), no. 2, 181--279, DOI 10.1016/S0012-9593(98)80015-3. MR1603257
\item[{[Fri82]}] Eric M. Friedlander, \emph{\'Etale homotopy of simplicial schemes}, Annals of Mathematics Studies, No. 104, Princeton University Press, Princeton, NJ; University of Tokyo Press, Tokyo, 1982. MR676809
\item[{[FW25]}] Tony Feng and Jonathan Wang, \emph{Geometric Langlands duality for periods}, Geom. Funct. Anal. \textbf{35} (2025), no. 2, 463--541, DOI 10.1007/s00039-025-00702-4. MR4880204
\item[{[FYZ21]}] Tony Feng, Zhiwei Yun, and Wei Zhang, \emph{Higher theta series for unitary groups over function fields} (English, with English and French summaries), Ann. Sci. \'Ec. Norm. Sup\'er. (4) \textbf{58} (2025), no. 2, 275--388. MR4902399
\item[{[FYZ23]}] Tony Feng, Zhiwei Yun, and Wei Zhang, \emph{Modularity of higher theta series I: cohomology of the generic fiber}, arXiv:2308.10979, 2023.
\item[{[FYZ24]}] Tony Feng, Zhiwei Yun, and Wei Zhang, \emph{Higher Siegel-Weil formula for unitary groups: the non-singular terms}, Invent. Math. \textbf{235} (2024), no. 2, 569--668, DOI 10.1007/s00222-023-01228-y. MR4689373
\item[{[Gai07]}] D. Gaitsgory, \emph{On de Jong's conjecture}, Israel J. Math. \textbf{157} (2007), 155--191, DOI 10.1007/s11856-006-0006-2. MR2342444
\item[{[Geh22]}] Lennart Gehrmann, \emph{Gelfand's trick for the spherical derived Hecke algebra}, Int. Math. Res. Not. IMRN \textbf{18} (2022), 13841--13856, DOI 10.1093/imrn/rnab130. MR4485946
\item[{[GJ09]}] Paul G. Goerss and John F. Jardine, \emph{Simplicial homotopy theory}, Modern Birkh\"auser Classics, Birkh\"auser Verlag, Basel, 2009. Reprint of the 1999 edition [MR1711612], DOI 10.1007/978-3-0346-0189-4. MR2840650
\item[{[Gro]}] Benedict Gross, \emph{Odd galois representations}, Available at https://people.math.harvard.edu/\textasciitilde gross/preprints/Galois\_Rep.pdf.
\item[{[GS07]}] Paul Goerss and Kristen Schemmerhorn, \emph{Model categories and simplicial methods}, Interactions between homotopy theory and algebra, Contemp. Math., vol. 436, Amer. Math. Soc., Providence, RI, 2007, pp. 3--49, DOI 10.1090/conm/436/08403. MR2355769
\item[{[GV18]}] S. Galatius and A. Venkatesh, \emph{Derived Galois deformation rings}, Adv. Math. \textbf{327} (2018), 470--623, DOI 10.1016/j.aim.2017.08.016. MR3762000
\item[{[Han17]}] David Hansen, \emph{Universal eigenvarieties, trianguline Galois representations, and $p$-adic Langlands functoriality}, with an appendix by James Newton, J. Reine Angew. Math. \textbf{730} (2017), 1--64, DOI 10.1515/crelle-2014-0130. MR3692014
\item[{[Har90]}] Michael Harris, \emph{Automorphic forms and the cohomology of vector bundles on Shimura varieties}, Automorphic forms, Shimura varieties, and $L$-functions, Vol. II (Ann Arbor, MI, 1988), Perspect. Math., vol. 11, Academic Press, Boston, MA, 1990, pp. 41--91. MR1044828
\item[{[Hid98]}] Haruzo Hida, \emph{Automorphic induction and Leopoldt type conjectures for $\mathrm{GL}(n)$}, Mikio Sato: a great Japanese mathematician of the twentieth century, Asian J. Math. \textbf{2} (1998), no. 4, 667--710, DOI 10.4310/AJM.1998.v2.n4.a5. MR1734126
\item[{[HLTT16]}] Michael Harris, Kai-Wen Lan, Richard Taylor, and Jack Thorne, \emph{On the rigid cohomology of certain Shimura varieties}, Res. Math. Sci. \textbf{3} (2016), Paper No. 37, 308, DOI 10.1186/s40687-016-0078-5. MR3565594
\item[{[HT17]}] David Hansen and Jack A. Thorne, \emph{On the $\mathrm{GL}_n$-eigenvariety and a conjecture of Venkatesh}, Selecta Math. (N.S.) \textbf{23} (2017), no. 2, 1205--1234, DOI 10.1007/s00029-017-0303-0. MR3624909
\item[{[HV19]}] Michael Harris and Akshay Venkatesh, \emph{Derived Hecke algebra for weight one forms}, Exp. Math. \textbf{28} (2019), no. 3, 342--361, DOI 10.1080/10586458.2017.1409144. MR3985839
\item[{[IKM22]}] Srikanth B. Iyengar, Chandrashekhar B. Khare, and Jeffrey Manning, \emph{Congruence modules and the Wiles-Lenstra-Diamond numerical criterion in higher codimensions}, Invent. Math. \textbf{238} (2024), no. 3, 769--864, DOI 10.1007/s00222-024-01292-y. MR4824729

\ysrc{566}
\item[{[Ill71]}] Luc Illusie, \emph{Complexe cotangent et d\'eformations. I} (French), Lecture Notes in Mathematics, Vol. 239, Springer-Verlag, Berlin-New York, 1971. MR491680
\item[{[Ill72]}] Luc Illusie, \emph{Complexe cotangent et d\'eformations. II}, Lecture Notes in Mathematics, Vol. 283, Springer-Verlag, Berlin-New York, 1972.
\item[{[Iye07]}] Srikanth Iyengar, \emph{Andr\'e--Quillen homology of commutative algebras}, Interactions between homotopy theory and algebra, Contemp. Math., vol. 436, Amer. Math. Soc., Providence, RI, 2007, pp. 203--234, DOI 10.1090/conm/436/08410. MR2355775
\item[{[Kha19]}] Adeel A. Khan, \emph{Virtual Fundamental Classes of derived stacks I}, Available at arXiv:1909.01332, 2019.
\item[{[Kis08]}] Mark Kisin, \emph{Potentially semi-stable deformation rings}, J. Amer. Math. Soc. \textbf{21} (2008), no. 2, 513--546, DOI 10.1090/S0894-0347-07-00576-0. MR2373358
\item[{[Koz16]}] Karol Kozio\l, \emph{A classification of the irreducible mod-$p$ representations of $U(1,1)(\Q_{p^2}/\Q_p)$} (English, with English and French summaries), Ann. Inst. Fourier (Grenoble) \textbf{66} (2016), no. 4, 1545--1582, DOI 10.5802/aif.3043. MR3494178
\item[{[KR14]}] Stephen Kudla and Michael Rapoport, \emph{Special cycles on unitary Shimura varieties II: Global theory}, J. Reine Angew. Math. \textbf{697} (2014), 91--157, DOI 10.1515/crelle-2012-0121. MR3281653
\item[{[KR23a]}] Chandrashekhar Khare and Niccol\`o Ronchetti, \emph{Derived Hecke action at $p$ and the ordinary $p$-adic cohomology of arithmetic manifolds}, Amer. J. Math. \textbf{145} (2023), no. 6, 1631--1694, DOI 10.1353/ajm.2023.a913294. MR4684289
\item[{[KR23b]}] Chandrashekhar Khare and Niccol\`o Ronchetti, \emph{Derived Hecke action at $p$ and the ordinary $p$-adic cohomology of arithmetic manifolds}, Amer. J. Math. \textbf{145} (2023), no. 6, 1631--1694, DOI 10.1353/ajm.2023.a913294. MR4684289\tn{底本中 [KR23a] 与 [KR23b] 两条文献内容完全相同（原书 p.566）。}
\item[{[KSZ21]}] Mark Kisin, Sug Woo Shin, and Yihang Zhu, \emph{The stable trace formula for Shimura varieties of abelian type}, 2021, to appear in Annals of Math. Studies.
\item[{[Kud04]}] Stephen S. Kudla, \emph{Special cycles and derivatives of Eisenstein series}, Heegner points and Rankin $L$-series, Math. Sci. Res. Inst. Publ., vol. 49, Cambridge Univ. Press, Cambridge, 2004, pp. 243--270, DOI 10.1017/CBO9780511756375.009. MR2083214
\item[{[KW09]}] Chandrashekhar Khare and Jean-Pierre Wintenberger, \emph{Serre's modularity conjecture. II}, Invent. Math. \textbf{178} (2009), no. 3, 505--586, DOI 10.1007/s00222-009-0206-6. MR2551764
\item[{[Li25]}] Chao Li, \emph{Geometric and arithmetic theta correspondences}, The Langlands program: III. Periods, $L$-functions, and the relative trace formula, Proc. Sympos. Pure Math., vol. 112.3, Amer. Math. Soc., Providence, RI, 2025.
\item[{[Lur04]}] Jacob Lurie, \emph{Derived algebraic geometry}, ProQuest LLC, Ann Arbor, MI, 2004. Thesis (Ph.D.)--Massachusetts Institute of Technology. MR2717174
\item[{[Lur09]}] Jacob Lurie, \emph{Higher topos theory}, Annals of Mathematics Studies, vol. 170, Princeton University Press, Princeton, NJ, 2009, DOI 10.1515/9781400830558. MR2522659
\item[{[Lur17]}] Jacob Lurie, \emph{Higher algebra}, 2017, Available at https://www.math.ias.edu/\textasciitilde lurie/papers/HA.pdf.
\item[{[Mac71]}] Saunders MacLane, \emph{Categories for the working mathematician}, Graduate Texts in Mathematics, Vol. 5, Springer-Verlag, New York-Berlin, 1971. MR354798
\item[{[Mad23]}] Keerthi Madapusi, \emph{Derived special cycles on Shimura varieties}, arXiv:2212.12849.
\item[{[Mar21]}] David Marcil, \emph{Numerical verification of a conjecture of Harris and Venkatesh}, J. Number Theory \textbf{221} (2021), 484--495, DOI 10.1016/j.jnt.2020.09.003. MR4203578
\item[{[Mata]}] Akhil Mathew, \emph{The Dold-Kan correspondence}, Available at http://math.uchicago.edu/\textasciitilde amathew/doldkan.pdf.
\item[{[Matb]}] Akhil Mathew, \emph{Is there a high-concept explanation for why ``simplicial'' leads to ``homotopy-theoretic''}, MathOverflow, https://mathoverflow.net/q/58497 (version: 2017-06-02).
\item[{[Matc]}] Akhil Mathew, \emph{Simplicial commutative rings, I}, Available at http://math.uchicago.edu/\textasciitilde amathew/SCR.pdf.
\item[{[May67]}] J. Peter May, \emph{Simplicial objects in algebraic topology}, Van Nostrand Mathematical Studies, No. 11, D. Van Nostrand Co., Inc., Princeton, N.J.-Toronto, Ont.-London, 1967. MR222892
\item[{[May92]}] J. Peter May, \emph{Simplicial objects in algebraic topology}, Chicago Lectures in Mathematics, University of Chicago Press, Chicago, IL, 1992. Reprint of the 1967 original. MR1206474

\ysrc{567}
\item[{[May99]}] J. P. May, \emph{A concise course in algebraic topology}, Chicago Lectures in Mathematics, University of Chicago Press, Chicago, IL, 1999. MR1702278
\item[{[Maz77]}] B. Mazur, \emph{Modular curves and the Eisenstein ideal}, with an appendix by Mazur and M. Rapoport, Inst. Hautes \'Etudes Sci. Publ. Math. \textbf{47} (1977), 33--186 (1978). MR488287
\item[{[Maz89]}] B. Mazur, \emph{Deforming Galois representations}, Galois groups over $\Q$ (Berkeley, CA, 1987), Math. Sci. Res. Inst. Publ., vol. 16, Springer, New York, 1989, pp. 385--437, DOI 10.1007/978-1-4613-9649-9\_7. MR1012172
\item[{[Mor25]}] Sophie Morel, \emph{Shimura varieties}, The Langlands program: I. Representations, Shimura varieties, and shtukas, Proc. Sympos. Pure Math., vol. 112.1, Amer. Math. Soc., Providence, RI, 2025.
\item[{[Ngo10]}] Bao Ch\^au Ng\^o, \emph{Le lemme fondamental pour les alg\`ebres de Lie} (French), Publ. Math. Inst. Hautes \'Etudes Sci. \textbf{111} (2010), 1--169, DOI 10.1007/s10240-010-0026-7. MR2653248
\item[{[Oll09]}] Rachel Ollivier, \emph{Le foncteur des invariants sous l'action du pro-$p$-Iwahori de $\mathrm{GL}_2(F)$} (French, with English summary), J. Reine Angew. Math. \textbf{635} (2009), 149--185, DOI 10.1515/CRELLE.2009.078. MR2572257
\item[{[OS19]}] Rachel Ollivier and Peter Schneider, \emph{The modular pro-$p$ Iwahori-Hecke Ext-algebra}, Representations of reductive groups, Proc. Sympos. Pure Math., vol. 101, Amer. Math. Soc., Providence, RI, 2019, pp. 255--308, DOI 10.1090/pspum/101/11. MR3930021
\item[{[PV21]}] Kartik Prasanna and Akshay Venkatesh, \emph{Automorphic cohomology, motivic cohomology, and the adjoint $L$-function} (English, with English and French summaries), Ast\'erisque \textbf{428} (2021), viii+132, DOI 10.24033/ast. MR4372499
\item[{[Qui]}] Daniel Quillen, \emph{Homology of commutative rings}, Available at chromotopy.org/paste/quillen.djvu.
\item[{[Qui67]}] Daniel G. Quillen, \emph{Homotopical algebra}, Lecture Notes in Mathematics, No. 43, Springer-Verlag, Berlin-New York, 1967. MR223432
\item[{[Qui70]}] Daniel Quillen, \emph{On the (co-) homology of commutative rings}, Applications of Categorical Algebra (Proc. Sympos. Pure Math., Vol. XVII, New York, 1968), Proc. Sympos. Pure Math., XVII, Amer. Math. Soc., Providence, RI, 1970, pp. 65--87. MR257068
\item[{[Ras25]}] Sam Raskin, \emph{An arithmetic application of Geometric Langlands}, The Langlands program: I. Representations, Shimura varieties, and shtukas, Proc. Sympos. Pure Math., vol. 112.1, Amer. Math. Soc., Providence, RI, 2025.
\item[{[Ron]}] Niccol\`o Ronchetti, \emph{A Satake homomorphism for the mod $p$ derived Hecke algebra}, arXiv:1808.06512.
\item[{[Sch12]}] Peter Scholze, \emph{Perfectoid spaces}, Publ. Math. Inst. Hautes \'Etudes Sci. \textbf{116} (2012), 245--313, DOI 10.1007/s10240-012-0042-x. MR3090258
\item[{[Sch15a]}] Peter Schneider, \emph{Smooth representations and Hecke modules in characteristic $p$}, Pacific J. Math. \textbf{279} (2015), no. 1-2, 447--464, DOI 10.2140/pjm.2015.279.447. MR3437786
\item[{[Sch15b]}] Peter Scholze, \emph{On torsion in the cohomology of locally symmetric varieties}, Ann. of Math. (2) \textbf{182} (2015), no. 3, 945--1066, DOI 10.4007/annals.2015.182.3.3. MR3418533
\item[{[Sch18]}] Peter Scholze, \emph{$p$-adic geometry}, Proceedings of the International Congress of Mathematicians---Rio de Janeiro 2018. Vol. I. Plenary lectures, World Sci. Publ., Hackensack, NJ, 2018, pp. 899--933.
\item[{[Sno18]}] Andrew Snowden, \emph{Singularities of ordinary deformation rings}, Math. Z. \textbf{288} (2018), no. 3-4, 759--781, DOI 10.1007/s00209-017-1912-1. MR3778977
\item[{[Sta71]}] H. M. Stark, \emph{Values of $L$-functions at $s=1$. I. $L$-functions for quadratic forms}, Advances in Math. \textbf{7} (1971), 301--343 (1971), DOI 10.1016/S0001-8708(71)80009-9. MR289429
\item[{[Sta75]}] H. M. Stark, \emph{$L$-functions at $s=1$. II. Artin $L$-functions with rational characters}, Advances in Math. \textbf{17} (1975), no. 1, 60--92, DOI 10.1016/0001-8708(75)90087-0. MR382194
\item[{[Sta76]}] H. M. Stark, \emph{$L$-functions at $s=1$. III. Totally real fields and Hilbert's twelfth problem}, Advances in Math. \textbf{22} (1976), no. 1, 64--84, DOI 10.1016/0001-8708(76)90138-9. MR437501

\ysrc{568}
\item[{[Sta80]}] Harold M. Stark, \emph{$L$-functions at $s=1$. IV. First derivatives at $s=0$}, Adv. in Math. \textbf{35} (1980), no. 3, 197--235, DOI 10.1016/0001-8708(80)90049-3. MR563924
\item[{[Sta22]}] The Stacks project authors, \emph{The stacks project}, https://stacks.math.columbia.edu, 2022.
\item[{[Su]}] Jun Su, \emph{Coherent cohomology of Shimura varieties and automorphic forms}, Ann. Sci. \'Ec. Norm. Sup\'er. (4) \textbf{57} (2024), no. 4, 1039--1099. MR4773300
\item[{[TV08]}] Bertrand To\"en and Gabriele Vezzosi, \emph{Homotopical algebraic geometry. II. Geometric stacks and applications}, Mem. Amer. Math. Soc. \textbf{193} (2008), no. 902, x+224, DOI 10.1090/memo/0902. MR2394633
\item[{[TV16]}] David Treumann and Akshay Venkatesh, \emph{Functoriality, Smith theory, and the Brauer homomorphism}, Ann. of Math. (2) \textbf{183} (2016), no. 1, 177--228, DOI 10.4007/annals.2016.183.1.4. MR3432583
\item[{[Urb11]}] Eric Urban, \emph{Eigenvarieties for reductive groups}, Ann. of Math. (2) \textbf{174} (2011), no. 3, 1685--1784, DOI 10.4007/annals.2011.174.3.7. MR2846490
\item[{[Vak24]}] Ravi Vakil, \emph{The rising sea---foundations of algebraic geometry}, Princeton University Press, Princeton, NJ, 2025. MR4942570
\item[{[Ven19]}] Akshay Venkatesh, \emph{Derived Hecke algebra and cohomology of arithmetic groups}, Forum Math. Pi \textbf{7} (2019), e7, 119, DOI 10.1017/fmp.2019.6. MR4061961
\item[{[Vez10]}] Gabriel Vezzosi, \emph{A note on the cotangent complex in derived algebraic geometry}, 2010, Available at https://arxiv.org/pdf/1008.0601.pdf.
\item[{[vS24]}] K\k{e}stutis \v{C}esnavi\v{c}ius and Peter Scholze, \emph{Purity for flat cohomology}, Ann. of Math. (2) \textbf{199} (2024), no. 1, 51--180, DOI 10.4007/annals.2024.199.1.2. MR4681144
\item[{[Yun18]}] Zhiwei Yun, \emph{Hitchin type moduli stacks in automorphic representation theory}, Proceedings of the International Congress of Mathematicians---Rio de Janeiro 2018. Vol. II. Invited lectures, World Sci. Publ., Hackensack, NJ, 2018, pp. 1447--1476. MR3966816
\item[{[Zhi25]}] Zhiwei Yun, \emph{Introduction to shtukas and their moduli}, The Langlands program: I. Representations, Shimura varieties, and shtukas, Proc. Sympos. Pure Math., vol. 112.1, Amer. Math. Soc., Providence, RI, 2025.
\item[{[Zha23]}] Robin Zhang, \emph{The Harris-Venkatesh Conjecture for Derived Hecke Operators}, ProQuest LLC, Ann Arbor, MI, 2023. Thesis (Ph.D.)--Columbia University. MR4593427
\item[{[Zhu25]}] Xinwen Zhu, \emph{Coherent sheaves on the stack of Langlands parameters}, The Langlands program: II. Geometrization of the Langlands correspondence, Proc. Sympos. Pure Math., vol. 112.2, Amer. Math. Soc., Providence, RI, 2025.
\end{refs}

\medskip
\begin{center}
\begin{minipage}{\linewidth}
\centering
{\small Department of Mathematics, University of California Berkeley, Berkeley, California 94720\\[0.3em]
\itshape Email address: \texttt{fengt@berkeley.edu}}

\medskip
{\small Department of Mathematics, Columbia University, New York, New York 10027\\[0.3em]
\itshape Email address: \texttt{harris@math.columbia.edu}}
\end{minipage}
\end{center}
